% BHU PYQ -- Manually transcribed & error-corrected
\documentclass[11pt,a4paper]{article}

\usepackage[a4paper,top=2.5cm,bottom=2.5cm,left=2.8cm,right=2.8cm,headheight=14pt]{geometry}
\usepackage{amsmath,amssymb}
\usepackage[T1]{fontenc}
\usepackage[utf8]{inputenc}
\usepackage{lmodern}
\usepackage{microtype}
\usepackage{fancyhdr}
\usepackage{enumitem}
\usepackage{hyperref}
\usepackage{lastpage}
\usepackage{parskip}

\hypersetup{colorlinks=true,urlcolor=black,linkcolor=black,
  pdftitle={CHB-04A: Ancillary Chemistry-II -- BHU PYQ},pdfauthor={Banaras Hindu University}}

\pagestyle{fancy}\fancyhf{}
\renewcommand{\headrulewidth}{0.4pt}\renewcommand{\footrulewidth}{0.4pt}
\fancyhead[L]{\small   \textit{Banaras Hindu University}}
\fancyhead[R]{\small   \textit{CHB-04A: Ancillary Chemistry-II}}
\fancyfoot[C]{\small Page  \thepage\ of \pageref{LastPage}}


\newlist{parts}{enumerate}{2}
\setlist[parts,1]{label=(\alph*),leftmargin=*,itemsep=5pt,topsep=3pt}
\setlist[parts,2]{label=(\roman*),leftmargin=*,itemsep=3pt,topsep=2pt}

\newcommand{\pts}[1]{\hfill   \textit{[#1]}}


\begin{document}


\begin{center}
  {\large\bfseries BANARAS HINDU UNIVERSITY}\\[2pt]
  {\normalsize B.Sc. (Hons.) Semester IV Examination, 2023-24}\\[6pt]
  \rule{0.8\linewidth}{0.5pt}\\[6pt]
  {\large\bfseries Paper No. CHB-04A: Ancillary Chemistry-II (Basic Aspects of Chemistry)}\\[4pt]
  {\normalsize Time: 3 Hours\hfill Maximum Marks: 70}
\end{center}
\rule{\linewidth}{0.4pt}
\small



\noindent   \textit{Instructions: Answer five questions in total, including Question~1 which is compulsory.}

\normalsize
\bigskip




\noindent   \textbf{Question 1.} Answer all parts of the following: \pts{5 $ \times$ 2 = 10}

\begin{parts}
  \item Write down the structure of Buna-S rubber and phenol-formaldehyde resin (Bakelite).
  \item Briefly discuss the different types of thermodynamic systems.
  \item Draw the Lewis dot structures for the following species: (i) $ \text{CO}_3^{2-}$ and (ii) $ \text{NH}_4 \text{Cl}$.
  \item Differentiate between binding energy and packing fraction.
  \item Define standard enthalpy of formation ($\Delta H_f^\circ$).
\end{parts}

\medskip\rule{\linewidth}{0.2pt}




\noindent   \textbf{Question 2.}

\begin{parts}
  \item Explain the significance of the octane number and cetane number in evaluating the quality of gasoline and diesel fuels, respectively. How do they correlate with engine performance? \pts{6}
  \item Discuss the potential health hazards associated with the production and use of synthetic polymers. How do polymer waste disposal and recycling contribute to environment preservation? \pts{8}
\end{parts}

\medskip\rule{\linewidth}{0.2pt}




\noindent   \textbf{Question 3.}

\begin{parts}
  \item Predict the sign of $\Delta G$, $\Delta H$, and $\Delta S$ for the sublimation of solid carbon dioxide ($ \text{CO}_2$) at room temperature. Give reasons. \pts{4}
  \item Assume a chemical reaction where $\Delta S$ is $+35.1\, \text{J/mol}\cdot \text{K}$ and $\Delta H$ is $-2.87\, \text{kJ/mol}$. Calculate the change in entropy for the universe ($\Delta S_{ \text{univ}}$) for this reaction at $60^\circ \text{C}$. Predict its spontaneity. \pts{6}
  \item State the second law of thermodynamics. \pts{4}
\end{parts}

\medskip\rule{\linewidth}{0.2pt}




\noindent   \textbf{Question 4.}

\begin{parts}
  \item Standard reduction potentials of the half-reactions are given below. Predict whether the following reaction will occur spontaneously or not. Give reasons:
    \[  \text{Fe}^{2+} +  \text{Cl}_2  o  \text{Fe}^{3+} + 2 \text{Cl}^- \]
    (Given: $E^\circ( \text{Fe}^{3+}/ \text{Fe}^{2+}) = +0.771\, \text{V}$, $E^\circ( \text{Cl}_2/ \text{Cl}^-) = +1.36\, \text{V}$). \pts{6}
  \item Standard reduction potentials ($E^\circ$) of some electrodes are: $ \text{A}^+/ \text{A} = -0.44\, \text{V}$, $ \text{B}^{2+}/ \text{B} = +0.34\, \text{V}$, and $ \text{C}^+/ \text{C} = +0.80\, \text{V}$. Predict the strongest oxidizing and strongest reducing agent among these. \pts{4}
  \item Describe the construction and working of the lead storage battery, showing charging and discharging cell reactions. \pts{4}
\end{parts}

\medskip\rule{\linewidth}{0.2pt}




\noindent   \textbf{Question 5.}

\begin{parts}
  \item Differentiate between Photosystem I and Photosystem II of photosynthesis in plants. \pts{6}
  \item Discuss the light and dark reactions of photosynthesis. \pts{5}
  \item State the biological role of haemoglobin and myoglobin. \pts{3}
\end{parts}

\vfill
\rule{\linewidth}{0.4pt}

\begin{center}\small
  \href{https://bhu-exam.vercel.app/}{Click here to visit our website}\\[4pt]
  \href{https://me-aryan.vercel.app/}{Click here to visit our contributor}\\[4pt]
  \href{https://quashan-me.vercel.app/}{Click here to visit our contributor}
\end{center}

\end{document}
